No início dos tempos, Donald E. Knuth criou o \TeX. Algum tempo depois, Leslie Lamport criou o \LaTeX\ \cite{lamport94latex} . Graças a eles, não somos obrigados a usar o Word nem o LibreOffice\footnote{Notas de rodapé também podem ser úteis.}.

A única intenção deste texto é mostrar como algumas coisas podem ser feitas. Existem muitos outros livros, manuais e tutoriais completos e muito melhores. Para o pessoal que quer começar a fazer e a apresentar os seus trabalhos de um jeito melhor, no entanto, pode começar por copiar este arquivo e editar sobre ele.

Algumas razões para você usar \LaTeX:

\begin{enumerate}
    \item Tabelas e figuras são colocadas no texto automaticamente. Você pode dar diretivas, ou pode deixar tudo com o \LaTeX. Não precisa ficar cuidando se vai ficar
    meia página em branco embaixo de uma figura. Em geral, as figuras e tabelas são colocadas o mais perto possível das referências a elas, mas no lugar que fica melhor,
    do ponto de vista estético.

    \item Referências cruzadas são feitas por meio de rótulos, ou \emph{labels}. Você pode rotular o que quiser, capítulos, seções, figuras, tabelas, equações. Rótulos
    são criados com o comando \verb"\label{nome-do-rotulo}". \label{nome-do-rotulo} Quando você quiser referenciar um rótulo, use o comando \verb"\ref{nome-do-rotulo}".
    Assim: o rótulo nome-do-rótulo está no item~\ref{nome-do-rotulo} desta lista e na página~\pageref{nome-do-rotulo} deste documento. Também está localizado na
    Seção~\ref{sec-figuras} do Capítulo~\ref{cap-introducao} deste texto. Referências às páginas são feitas com o comando \verb"\pageref{nome-do-rotulo}".

    \item Todas as numerações (de capítulos, seções, subseções, figuras, equações e páginas) são geradas automaticamente. Assim, você não precisa se preocupar, caso
    troque uma figura de lugar ou passe uma seção inteira para um outro capítulo, de ter que verificar se as numerações e referências estão ou não corretas.

    \item A bibliografia, assim como todas as citações bibliográficas, também podem ser gerada automaticamente.
\end{enumerate}


Em \LaTeX, o número de espaços em branco no arquivo fonte não conta. Assim, colocar uma quantidade maior de espaços entre os caracteres não
influi no resultado final.

Parágrafos são criados por meio de uma linha em branco no arquivo \texttt{.tex}.

O formato final (gerado como um arquivo com extensão \texttt{.dvi}) se for usado o compilador \texttt{\small latex}, é determinado pelo arquivo de formatação usado, indicado com o comando \texttt{$\backslash$documentclass}. Este arquivo, por exemplo, usa o formato \emph{destufba}, que formata trabalhos acadêmicos. O compilador correto para usar neste caso é o \texttt{\small pdflatex}, que gera os arquivos diretamente em formato PDF (\emph{portable document format}). Outros podem ser construídos ou usados. O texto não precisa ser alterado, somente o arquivo de formatação usado.


\section{Figuras e tabelas}
\label{sec-figuras}

Esta sessão faz referência às \figurename s~\ref{fig:estrutura},~\ref{fig:ex1} e~\ref{fig:ex2}, a título de exemplo. A primeira figura apresenta a estrutura de uma figura. A \emph{descrição} (ou \emph{caption}) deve aparecer \textbf{acima} da figura. Abaixo da figura, deve ser indicado a origem da imagem, mesmo se essa for apenas os autores do texto. Note que ambos os ambientes são ``flutuantes''. Ou seja, são colocados no texto na melhor posição possível, a partir das dicas fornecidas (ou não) pelo usuário. Veja que ao lado do início do ambiente encontra-se o texto \texttt{[htbp]}. \texttt{h} significa \emph{here} e indica que a figura deve, preferencialmente, ser colocada no ponto do texto onde a definição aparece; \texttt{t} significa \emph{top} e indica que a figura deve ser colocada no topo da página; \texttt{b} significa \emph{bottom} e indica que a figura deve ser colocada ao final da página; e \texttt{p} significa \emph{page} e indica que a figura deve aparecer sozinha em outra página. Usualmente o arquivo de formatação define como as figuras devem ser colocadas. Use somente se você não gostar do resultado, mas lembre que o \LaTeX\ entende mais de formatação do que você...

Note como as referências cruzadas são feitas em \LaTeX: cada estrutura a ser referenciada (figuras, tabelas, equações, itens de listas, capítulos, seções, etc.) pode ser rotulada com o comando $\backslash$label\{nome do rótulo\}, em qualquer ambiente numerado. O número capturado é o do ambiente numerado mais próximo (neste caso aqui, seria o número da \sectionname\ Figuras e Tabelas). Quando o rótulo está dentro de um ambiente figure, por exemplo, o número é o nome da figura gerada por este ambiente. Toda a numeração é gerada automaticamente pelo compilador: se você mudar uma figura de lugar, os cálculos serão refeitos e toda a numeração será gerada corretamente. Compile o texto duas vezes seguidas para que todas as referências sejam calculadas corretamente.

A \figurename~\ref{fig:ex1} representa o caso mais comum, onde a figura propriamente dita é importada de um arquivo (neste exemplo em formato \texttt{eps} ou \texttt{pdf}. Veja a \sectionname~\ref{sec:fig_format}). A \figurename~\ref{fig:ex2} exemplifica o uso do environment \texttt{picture}, para desenhar usando o próprio~\LaTeX\ \cite{mittelbach04thelatex, goossens04thelatexgraphics,lemke02zen}.

\begin{figure}[htbp]
    \caption{Descrição da Figura deve ir no topo}
    \begin{center}
        % Aqui vai um includegraphics , um picture environment ou qualquer
        % outro comando necessário para incorporar o formato de imagem
        % utilizado.
        \begin{picture}(100,100)
            \put(0,0){\line(0,1){100}}
            \put(0,0){\line(1,0){100}}
            \put(100,100){\line(0,-1){100}}
            \put(100,100){\line(-1,0){100}}
            \put(10,50){Uma Imagem}
        \end{picture}
    \end{center}
    \label{fig:estrutura}
    \fonte{Os autores}
\end{figure}

\begin{figure}[b]
    \caption{Exemplo de figura importada de um arquivo e também exemplo de caption muito grande que ocupa mais de uma linha na Lista~de~Figuras}
    \centerline{\includegraphics[width=10em]{Dest_figure.png}}
    \fonte{https://
\end{figure}

\begin{figure}[p]
    \caption{Exemplo de figura desenhada com o environment \texttt{picture}, que aparece sozinha numa página.}
    \begin{center}
        \setlength{\unitlength}{.1em}
        \begin{picture}(100,100)
            \put(20,20){\circle{20}}
            \put(20,20){\small\makebox(0,0){a}}
            \put(80,80){\circle{20}}
            \put(80,80){\small\makebox(0,0){b}}
            \put(28,28){\vector(1,1){44}}
        \end{picture}
    \end{center}
    \fonte{Os autores}
    \label{fig:ex2}
\end{figure}

Tabelas são construídas com praticamente os mesmos comandos. Ver as \tablename s~\ref{tab-isto-e-uma-tabela} e \ref{tab-tabela-exemplo}. No exemplo as tabelas estão com posicionamento forçado, mas é melhor não usar. Deixar com o \LaTeX\ é sempre uma melhor política, esteticamente falando.

\begin{table}[h]
    \caption{Isto é uma tabela}
    \centering
    \begin{tabular}{|r|c|l|}
        % after \\: \hline or \cline{col1-col2} \cline{col3-col4} ...
        \hline
        Texto         & Texto         & Texto         \\ \hline
        este texto    & este texto    & este texto    \\
        é posicionado & é posicionado & é posicionado \\
        à direita     & no centro     & à esquerda    \\ \hline
    \end{tabular}
    \fonte{A autora}
    \label{tab-isto-e-uma-tabela}
\end{table}


\begin{table}[t]
    \caption{Uma tabela de Exemplo}
    \begin{center}
        \begin{tabular}{c|c|p{5cm}}
            \textit{Col 1} & \textit{Col 2} & \textit{Col 3}                                                                                                                                \\
            \hline
            \hline
            Val 1          & Val 2          & Esta coluna funciona como um parágrafo, tendo uma margem definida em 5cm. Quebras de linha funcionam como em qualquer parágrafo do documento. \\
            Valor Longo    & Val 2          & Val 3                                                                                                                                         \\
            \hline
        \end{tabular}
    \end{center}
    \fonte{Os autores}
    \label{tab-tabela-exemplo}
\end{table}

\subsection{Formato de Figuras}
\label{sec:fig_format}

O \LaTeX permite utilizar vários formatos de figuras, entre eles \emph{eps}, \emph{pdf}, \emph{jpeg} e \emph{png}. Programas de diagramação como Inkscape (e mesmo LibreOffice) permitem gerar arquivos de imagens vetoriais que podem ser utilizados pelo LaTeX sem dificuldade. Pacotes externos permitem utilizar SVG e outros formatos.

\subsection{Classificação dos etc.}

O formato adotado pela ABNT prevê apenas três níveis (capítulo, seção e subseção). Assim, \texttt{subsubsection} não deve ser usado. Não só pela recomendação como pela feiúra do estilo. Aliás, evite usar subseções. A quebra em níveis torna a leitura menos fluente e dois níveis, com raríssimas exceções, são suficientes para trabalhos acadêmicos em geral.

Citações têm um formato específico também. Use quando a citação possuir mais do que 3 linhas. Em citações é sempre obrigatório colocar a fonte e o número da página (e não como está aqui...).

\begin{quote}
    Uma citação é um texto copiado ipsis literis de algum outro lugar. Usa-se somente quando se quer apresentar a ideia do autor exatamente como foi produzida, para analisá-la em detalhe.
\end{quote}


\section{Equações matemáticas}

No corpo do texto para escreve-se equações e expressões matemáticas entre dois símbolos \$. Subscritos são feitos com um \emph{underscore} e superscritos com um acento circunflexo. Por exemplo faz-se um $x^2$ com a expressão \verb"$x^2$".

Outras maneiras de se fazer equações são com os comandos \verb"\[" e \verb"\]", que geram uma expressão centralizada. Por exemplo, a expressão \verb"\[\Delta=b^2-4ac\]" gera a seguinte saída:

\[\Delta=b^2-4ac\]

Expressões mais complicadas também podem ser feitas, em modo texto, e geram saídas bonitinhas. Por exemplo:

\[\Phi(x,y,z)=\sqrt[3]{\frac{x^2+y^2+z^2}{(x+y+z)^{1/3}}}\]

Matrizes também podem ser geradas facilmente:

\[M_{m\times n}=\left( \begin{array}{cccc}
                           a_{11} & a_{12} & \ldots & a_{1n} \\
                           a_{21} & a_{22} & \ldots & a_{2n} \\
                           \ldots & \ldots & \ldots & \ldots \\
                           a_{m1} & a_{m2} & \ldots & a_{mn} \\
\end{array}\right)\]

Se você precisar numerar equações ou expressões, use o ambiente equation:

\begin{equation}
    \label{minha-equacao}
    \Theta(\overrightarrow{x})=\frac{\partial\nabla^2}{\partial t}
\end{equation}

Ele faz o mesmo que \verb"\[\]", mas gera um número que permite
referências cruzadas. Por exemplo, a Equação~\ref{minha-equacao} é
bem fácil de fazer, e a Equação~\ref{outra-equacao} também.

\begin{equation}
    \label{outra-equacao}
    \Theta(\overrightarrow{x})=
    \left\{\begin{array}{lcl}
               \frac{x^2+y^2}{x-y} & & x\geq y \\
               0                   & & x = y   \\
               \frac{x^2+y^2}{y-x} & & x\leq y
    \end{array}\right.
\end{equation}


\section{Sobre as referências bibliográficas}

A classe \emph{destufba} faz uso do pacote \emph{abnTeX2} com algumas alterações
feitas por Sandro Rama Fiorini. Culpe ele se algo der errado. Agradeça a ele
pelo que der certo. As modificações dão uma camada de tinta NATBIB-style,
já que o abntex2 usa uns comandos de citação feitos para alienígenas de 5 braços. Exemplos de citação diferentes para a mesma obra (\cite{abramski92handbook1}):

\begin{itemize}
    \item \emph{cite}: Unicórnios são verdes \cite{abramski92handbook1};
    \item \emph{citep}:Unicórnios são verdes \citep{abramski92handbook1};
    \item \emph{citet}: Segundo \citet{abramski92handbook1}, unicórnios são
    verdes.
    \item \emph{citen or citenum}: Segundo \citen{abramski92handbook1},
    unicórnios são verdes.
    \item \emph{citeauthor e citeyearpar}: Segundo artigos de
    \citeauthor{abramski92handbook1}, unicórnios são verdes
    \citeyearpar{abramski92handbook1}.
\end{itemize}

As referências bibliográficas devem ser todas colocadas em um arquivo .bib. Junto com esse pacote também vai um exemplo. Somente as referências citadas no texto com o comando $\backslash$cite ou forçadas a aparecer com o comando $\backslash$nocite aparecerão na lista de referências produzida. Isso permite que você tenha um único arquivo de referências bibliográficas que pode ser usado em qualquer trabalho e incrementado ao longo do tempo. As referências são usadas para qualquer modelo de bibliografia, então podem ser usadas para artigos, apresentações, etc.

